\documentclass{article}

\usepackage{listings}
\usepackage[utf8]{inputenc}
\usepackage[spanish]{babel} % Agregar es-nodecimaldot si hay problemas con el separador de decimales y de miles
% \usepackage{noto-mono} % Tipografía "Noto Sans Mono" para monospace
\usepackage{amsmath, amsbsy, amssymb} % Agregar cuando se necesiten: amscd, amssymb, amsthm, latexsym
\usepackage{enumerate}
\usepackage{graphicx}
\usepackage{multicol}
\usepackage{svg}
\usepackage{tabularx}
\usepackage{algorithm}
\usepackage{algpseudocode}
\usepackage{caption}
\usepackage{etoolbox}

% Tamaño de hoja y márgenes
\usepackage[a4paper,top=1cm,bottom=2cm,left=1cm,right=1cm]{geometry}
\setlength{\parindent}{0em}
\setlength{\parskip}{1em}
\setlength{\intextsep}{1em}
\setlength{\abovedisplayskip}{1em}
\setlength{\belowdisplayskip}{1em}
\setlength{\abovedisplayshortskip}{1em}
\setlength{\belowdisplayshortskip}{1em}

% Numeración de las secciones y el índice
\setcounter{tocdepth}{1}
\renewcommand{\thesubsection}{\thesection.\alph{subsection}}

% Color de los links externos
\usepackage{hyperref}
\hypersetup{colorlinks=true, linkcolor=black, urlcolor=blue}

% Permite math mode adentro de listings (bloques de "código")
\usepackage{listings}
\lstset{
    inputencoding=utf8,
    extendedchars=\true,
    basicstyle=\ttfamily\small,
    mathescape,
    literate=%
        {á}{{\'a}}1
        {é}{{\'e}}1
        {í}{{\'i}}1
        {ó}{{\'o}}1
        {ú}{{\'u}}1
        {Á}{{\'A}}1
        {É}{{\'E}}1
        {Í}{{\'I}}1
        {Ó}{{\'O}}1
        {Ú}{{\'U}}1
        {ñ}{{\~n}}1
        {Ñ}{{\~N}}1
}

% Alias para íconos/símbolos
\usepackage{pifont}
\newcommand{\xmark}{\color{purple}\ding{54}}

% Otros alias
\newcommand{\xor}{\oplus}
\newcommand{\nor}{\downarrow}
\newcommand{\opsub}[2]{\ensuremath{#1_{\mathrm{#2}}}}
\newcommand{\yLuego}{\opsub{\land}{\scriptscriptstyle{L}}}
\newcommand{\oLuego}{\opsub{\lor}{\scriptscriptstyle{L}}}
\newcommand{\implicaLuego}{\opsub{\implies}{\scriptscriptstyle{L}}}
\renewcommand{\implies}{\Rightarrow}

% Columna "x" tiene fondo amarillo
\usepackage{colortbl}
\newcolumntype{x}{>{\columncolor[HTML]{FFF2CC}}c}

% Tipografía para los algoritmos
\AtBeginEnvironment{algorithmic}{\small}

% Saca sufijo "Algorithm X"
\captionsetup[algorithm]{labelformat=empty}

% Alias para algoritmos
\newcommand{\Complexity}[1]{\textbf{Complejidad}: #1}
\newcommand{\Pre}[1]{\textbf{Pre} $\equiv$ \{#1\}}

% Carátula de prácticas
%practicas
\newcommand{\practica}[2]{%
    \title{Pr\'actica #1 \\ #2}
    \author{Algoritmos y Estructuras de Datos II}
    \date{Segundo Cuatrimestre 2022}

    \maketitlepractica{#1}{#2}
}

\newcommand \maketitlepractica[2] {%
\begin{center}
\begin{tabular}{r cr}
 \begin{tabular}{c}
{\large\bf\textsf{\ Algoritmos y Estructuras de Datos II\ }}\\
Segundo Cuatrimestre 2022\\ \\
Práctica #1 - \textbf{#2}\\
\end{tabular} &
\begin{tabular}{@{} p{1.6cm} @{}}
\includegraphics[width=4.5cm]{./logodtp.png}
\end{tabular} &
\begin{tabular}{l @{}}
 \emph{Departamento de Computaci\'on} \\
 \emph{Facultad de Ciencias Exactas y Naturales} \\
 \emph{Universidad de Buenos Aires} \\
\end{tabular}
\end{tabular}
\end{center}

\bigskip
}
\setlength{\parindent}{0pt}

\title{Modelo Primer Parcial}
\author{TADs y especificación}
\date{Septiembre 2022}

\begin{document}

\maketitle

\section{Ejercicio 1}

tieneValor: $dicc(\alpha, nat) \times nat \rightarrow bool$

tieneValor(d, n) $\equiv $
\begin{lstlisting}
    if vacio?(claves(d))
        then false
        else if obtener(d, dameUno(claves(d))) = n
            then true
            else tieneValor(borrar(DameUno(claves(d)), d), n)
\end{lstlisting}
\section{Ejercicio 2}

diasRestantesDeAlquiler(inaugurarVC(cs,ps), c2) $\equiv$ vacio

diasRestantesDeAlquiler(alquilar(v,c,p), c2) $\equiv$ 
\begin{lstlisting}
    if c = c2
        then definir(p, obtener(p, pelis(v)), diasRestantesDeAlquiler(v,c))
        else diasRestantesDeAlquiler(v, c2)
\end{lstlisting}

diasRestantesDeAlquiler(devolver(v,c,p), c2) $\equiv$ 
\begin{lstlisting}
    if c = c2
        then borrar(p, diasRestantesDeAlquiler(v,c))
        else diasRestantesDeAlquiler(v,c)
\end{lstlisting}

diasRestantesDeAlquiler(pasarDia(v), c2) $\equiv$ restarUno(quitarVencidas(diasRestantesDeAlquiler(v, c2)))

clientes(inaugurarVC(cs, ps)) $\equiv$ cs

clientes(alquilar(v,c,p)) $\equiv$ clientes(v)

clientes(devolver(v,c,p)) $\equiv$ clientes(v)

clientes(pasarDia(v)) $\equiv$ clientesSinVencidas(c, clientes(v))

quitarVencidas(alqs) $\equiv$ borrarLasVencidas(alqs, claves(alqs))

borrarLasVencidas(d, c) $\equiv$ 
\begin{lstlisting}
    if vacio?(c)
        then vacio
        else if obtener(dameUno(c), d) = 0
            then borrarLasVencidas(borrar(dameUno(c), d), c)
            else definir(dameUno(c), obtener(dameUno(c), d), borrarLasVencidas(d, sinUno(c)))
\end{lstlisting}

clientesSinVencidas(v, cs) $\equiv$ 
\begin{lstlisting}[mathescape=true]
    if vacio?(cs)
        then $\emptyset$
        else if tieneValor(diasRestantesDeAlquiler(v, dameUno(cs)),0)
            then clientesSinVencidas(v, sinUno(cs))
            else Ag(dameUno(cs), clientesSinVencidas(v, sinUno(cs))
\end{lstlisting}

\section{Ejercicio 3}

\subsection*{Pregunta a}
Dos juegos de calamares se pueden diferenciar observando:
\begin{itemize}
    \item Los jugadores vivos
    \item Los seguidores que tiene cada jugador
\end{itemize}
    
\subsection*{Pregunta b}

TAD CALAMAR es String

\textbf{Observadores básicos}
\begin{itemize}
    \item jugadores: jc $\rightarrow$ conj(calamar)
    \item seguidores: calamar c $\times$ jc j $\rightarrow$ conj(calamar) \{c $\in$ jugadores(j)\}
\end{itemize}


\textbf{Igualdad observacional}
\begin{align*}
(\forall j, j': jc)(j =_{obs} j' \iff &jugadores(j) =_{obs} jugadores(j') \\
                    \wedge &(\forall c: calamar)(c \in jugadores(j) \implies_{L} seguidores(c,j) = seguidores(c, j'))
\end{align*}
   
\subsection*{Pregunta c}

\textbf{Generadores}
\begin{itemize}
    \item iniciarJuego: conj(calamates) $\rightarrow$ jc
    \item seguir: calamar a $\times$ calamar b $\times$ jc j $\rightarrow$ jc 
    
    \{a $\in$ jugadores(j) $\wedge$ b $\in$ jugadores(j) $\wedge$ $\neg$ esSeguidor(a, j)\}
    \item luchar: calamar a $\times$ calamar b $\times$ jc j $\rightarrow$ jc 
    
    \{a $\in$ jugadores(j) $\wedge$ b $\in$ jugadores(j)\}
\end{itemize}

\subsection*{Pregunda d}

jugadores(iniciarJuego(cs)) $\equiv$ cs

jugadores(seguir(a, b, j)) $\equiv$ jugadores(j)

jugadores(luchar(a, b, j)) $\equiv$ jugadores(j) - b

seguidores(c, iniciarJuego(cs)) $\equiv \emptyset$

seguidores(c, seguir(a,b,j)) $\equiv$

\begin{lstlisting}
    if c = b
        then Ag(a, seguidores(c,j))
        else seguidores(c,j)
\end{lstlisting}

seguidores(c, luchar(a,b,j)) $\equiv$
\begin{lstlisting}[mathescape=true]
    if c = a
        then seguidores(b, j) $\cup$ seguidores(c, j)
        else if c = b
            then $\emptyset$
            else seguidores(c,j)
\end{lstlisting}

\end{document}